\setscript[hanzi]
\setupalign[hanging, hz]

\usemodule[zhfontsext]
\usetypescript[cnfontb]
\setupbodyfont[cnfontb,rm,12pt] % 小四号 = 12pt
\mainlanguage[cn]

\setupinterlinespace[line=4.575ex] %line=3.05ex相当于WORD 中的1.0倍行距, 4.575ex相当于WORD 中的1.5倍行距

\setuppapersize[A4]
\setuplayout[width=fit,height=fit,backspace=10mm,cutspace=10mm,topspace=10mm,margin=10mm,header=10mm,footer=10mm]
\setuppagenumbering[location={header,right},style=\bfb]


\starttext

\midaligned{
三字经\\
\underbar{人之初，性本善。\underbar{性相近，习相远。}}\quad
\overbar{苟不教，性乃迁。\overbar{教之道，贵以专。}}\\
昔孟母，择邻处。子不学，断机杼。\quad
窦燕山，有义方。教五子，名俱扬。\\
\understrike{养不教，父之过。教不严，师之惰。}\quad
\overstrike{子不学，非所宜。 \overstrike{幼不学，老何为。}}\\
玉不琢，不成器。人不学，不知义。\quad
为人子，方少时。亲师友，习礼仪。\\
\underdot{香九龄，能温席。\underdot{孝于亲，所当执。}}\quad
\underdash{融四岁，能让梨，\underdash{悌于长，宜先知。}}\\
首孝悌，次见闻。知某数，识某文。\quad
一而十，十而百。百而千，千而万。\\
\underbar{\underdot{三才者，天地人。三光者，日月星。}}\quad
三纲者，君臣义。父子亲，夫妇顺。\\
}

\stoptext